\section{Research Questions}
\label{sec:intro:rqs}
Throughout, \emph{optimizability} is the fraction of AIG nodes a given synthesis algorithm
removes from a circuit (formalised in~\ref{sec:prelim:formalization}
and~\ref{sec:method:data:label:definition}), and \emph{compression ratio} is the node or edge
retention of a reduced graph relative to its source
(\ref{sec:prelim:reduction:measuring}). The distinction between node and edge retention is
not cosmetic: some reduction methods preserve node counts exactly and only remove edges, so
a single ``compression'' figure is ambiguous and RQ4 depends on the definition being pinned
down first.

\subsection{RQ1: Task Feasibility and Baseline}
\label{sec:intro:rq1}
How well can a Graph Neural Network predict the optimizability of a circuit directly from
its full, unreduced AIG representation, relative to trivial predictors, standard graph-level
GNN encoders, and published circuit-representation models?

\paragraph{RQ1a (protocol sensitivity).}
How much of the measured score is a property of the evaluation protocol rather than of the
model? Three train/test protocols are run under otherwise identical conditions: a plain
random split, a recipe-disjoint split, and the design-disjoint split used for every reported
result (\ref{sec:method:experiment:splitting}). The gap between them is a direct
measurement of how much of the apparent accuracy under the more permissive protocols is
design recognition rather than structural inference. The per-design breakdown of the
design-disjoint test set then shows whether the remaining error is uniform across held-out
designs or concentrated in a few (\ref{sec:results:rq1:protocol}).

\subsection{RQ2: Reduction Efficiency}
\label{sec:intro:rq2}
How do graph reduction techniques, spanning partitioning, sparsification, and
summarization/coarsening, differ in graph shrinkage rate, peak memory usage,
training and inference speed, and the offline cost of applying the reduction itself?

\subsection{RQ3: Predictive Retention and Trade-offs}
\label{sec:intro:rq3}
To what extent do models trained on reduced AIGs retain predictive accuracy (measured via
Smooth L1 Loss, RMSE, and $R^2$) and circuit-ranking ability (Spearman correlation)
relative to the full-graph baseline established in RQ1, and which techniques offer the
best accuracy-to-efficiency trade-off?

\paragraph{Hypothesis H1 (path contraction).}
Sparsification \emph{removes edges} and can only impede message propagation, whereas
summarization \emph{contracts paths} and shortens the distance a signal must travel. AIGs
are very deep relative to the encoder's four message-passing layers, so the model is
receptive-field starved and subject to over-squashing \citep{alon2021oversquashing}.
Coarsening therefore has a mechanism by which it could \emph{raise} accuracy rather than
trade it away, and is expected to retain more than sparsification at matched compression.
% TODO: the receptive-field metric is specified in
% \ref{sec:method:experiment:metrics:reduction} but not implemented, so H1 is
% currently asserted rather than tested.

\subsection{RQ4: Value of Domain-Informed Adaptation}
\label{sec:intro:rq4}
Do AIG-specific reduction heuristics, which exploit topological level, reconvergence
structure, and gate roles, retain more predictive accuracy than generic techniques at
equivalent compression? Where no generic method compresses to the same degree, the
comparison is made against a compression-matched random control. If domain knowledge does
not pay, what does that indicate about which structural features the model actually relies
on?

\paragraph{Hypothesis H2 (weak heuristics).}
The domain-informed adaptations tested here are relatively simple (edge-span weighting,
level slicing, gate-role filtering, reconvergence-bounded merging) and are not expected to
win outright. A negative result is still a result: it would indicate that the predictive
signal is distributed across the AIG rather than concentrated in the causal cones and long
paths these heuristics were built to protect. Section~\ref{sec:results:rq4:adequacy} bounds what
such a result would and would not license.

\subsection{RQ5: Reduced to Unreduced? Generalization}
\label{sec:intro:rq5}
Can a model trained exclusively on reduced AIGs (partitioned, sparsified, or summarized) be
queried on normal, full-sized AIGs at inference time, and does it retain predictive accuracy
across this change of structural state?

Inference is the regime where this pays off. Training stores activations and gradients for
every node in a batch, whereas forward-only inference stores neither. A model that had to be
trained on reduced graphs to fit in memory may therefore still serve full graphs on
substantially more modest hardware. The practical claim under test is \emph{train cheap,
infer full}.

% Positive control: colour refinement at unbounded count cap is provably lossless for
% this encoder (\ref{sec:method:architecture:exact}), so without it a poor cross-state
% result cannot be told apart from a broken evaluation path.

\subsection{Scope and Delimitations}
\label{sec:intro:rqs:scope}
Each item below states what is included and what is deliberately excluded. These are
choices made in advance; limitations discovered during the work are separated out
into~\ref{sec:discussion:limitations}.

\subsubsection{Target Synthesis Algorithm}
\label{sec:intro:scope:algorithm}
Training and evaluation target a single synthesis script, \texttt{Orchestrate}. The dataset
also contains graphs optimized by \texttt{Deepsyn}, \texttt{Syn4} and \texttt{C2RS}, and
those are left unused here. Focusing on one algorithm removes cross-algorithm variance as a
confound: the question is how much a \emph{reduction} costs, and comparing reductions across
several target algorithms at once would not isolate that. Extending to the remaining three
is the most direct piece of future work (\ref{sec:conclusion:futurework:algorithms}).

\subsubsection{Definition of Optimizability}
\label{sec:intro:scope:target}
Optimizability is defined as relative \emph{node} reduction only. Area, delay and power,
the objectives synthesis actually optimizes for, are not modelled, and node count is used
as a proxy for area. Whether that proxy is adequate is a construct-validity question taken
up in~\ref{sec:discussion:limitations:validity}.

\subsubsection{Graph Scale and Dataset Boundaries}
\label{sec:intro:scope:scale}
The corpus is bounded at an intermediate scale (see~\ref{sec:method:experiment:scale}), not
at industrial scale. This is a deliberate consequence of the design: RQ1 requires a
full-graph baseline, and a baseline that does not fit in memory cannot be measured. The
motivation invokes AIGs of millions to billions of nodes; the evaluation does not reach
them, and~\ref{sec:discussion:limitations:scalability} states plainly what does and does
not extrapolate.

\subsubsection{Reduction Families Considered}
\label{sec:intro:scope:reductions}
Three families are covered: partitioning, sparsification, and summarization/coarsening.
Graph \emph{condensation} (synthesising a small graph by gradient or distribution matching)
is excluded, for reasons given in~\ref{sec:relwork:reduction:summarization}: it produces no
correspondence between synthetic and real nodes, it is label-dependent by construction, and
it is tied to a specific architecture, so it cannot answer RQ5, which requires the same
weights to ingest both reduced and full graphs.

\subsubsection{Model Architecture}
\label{sec:intro:scope:architecture}
One encoder family is used throughout (an edge-aware GCN+,~\ref{sec:method:architecture}),
held fixed so that differences in outcome are attributable to the reduction rather than to
the model. Published circuit models are used as \emph{baselines} for RQ1
(\ref{sec:method:architecture:baselines}) but are not themselves trained on reduced graphs.
Whether the reduction rankings transfer to other encoders is untested.

\subsubsection{Prediction Rather Than Script Generation}
\label{sec:intro:scope:prediction}
This work predicts optimizability; it does not generate synthesis scripts. The gap between
the two is worth stating precisely, because the motivation is script selection and the
metrics are not. Spearman correlation here ranks \emph{circuits} by how optimizable they
are under one fixed script. Selecting a script would require ranking \emph{algorithms} for
one circuit, which a single-algorithm model cannot do. Closing that loop is future work
(\ref{sec:conclusion:futurework:generation}); the contribution claimed here is the
representation and reduction machinery it would need, not the selector itself.
